


\documentclass{article}
\usepackage{graphicx} % Required for inserting images
\usepackage[a4paper, total={6in, 8in}]{geometry}

\title{Optimization of a UAV for Wildfire Management}
\author{Nathaniel Hargan}
\usepackage{amsmath}
\date{}

\begin{document}
	
	\maketitle
	
	\large
	
	The problem is formulated as follows:
	$$ \min \textnormal{mass}(d_a,t_a,d_s,t_s,r_s,t_f,t_w) $$
	
	subject to:
	
	$$g_j(x) \le 0; j = 0, 1, 2, ..., J$$
	$$x_i^L \le x_i \le x_i^U$$
	
	design variables: 
	
	\begin{itemize}
		\item $d_a$: arm diameter,
		\item $t_a$: arm thickness,
		\item $d_s$: strut diameter,
		\item $t_s$: strut thickness,
		\item $r_s$: strut distance from center,
		\item $t_f$: hub flange thickness,
		\item $t_w$: hub web thickness 
	\end{itemize}
	
	derived variables:
	\begin{itemize}
		\item $d_{tip}$: Deflection at tip
		\item $\sigma_{max}$: Maximum stress from bending 
		\item $f_{amp}$: Amplitude of displacement at the tip of the uav arm
		\item $D$: Damage per mission
		\item $ E_{used}$: Energy used over th duration of the mission
		\item $ E_{total}$: Energy used over th duration of the mission
		
	\end{itemize}
	
	design parameters:
	\begin{itemize}
		\item $d_{allowable}$: Deflection allowable
		\item $\sigma_{allowable}$: Stress allowable from bending 
		\item $R_{allowable}$: Allowable amplitude displacement ratio
		\item $n_{missions} $: numbr of missions before damage
	\end{itemize}
		
	Constraints: \\
	\large
	
	Limit on deflection from bending by force required to lift the UAV. \\
	$$ g_1(d_a,t_a,d_s,t_s,r_s,t_f,t_w) = d_{tip} \le d_{allowable} $$
	
	Limit on stress from bending by force required to lift the UAV. \\
	$$ g_2(d_a,t_a,d_s,t_s,r_s,t_f,t_w) = \sigma_{max} \le \sigma_{allowable} $$
	
	
	Limit on amplitude of displacement given a forcing function based on the propeller speed of the UAV. \\
	$$ g_3(d_a,t_a,d_s,t_s,r_s,t_f,t_w) = \frac{f_{amp}}{d_{tip}} \le R_{allowable} $$
	
	
	Limit on the damage the UAV during a mission. \\
	$$ g_4(d_a,t_a,d_s,t_s,r_s,t_f,t_w) = D * n_{missions} \le 1 $$
	
	Energy
	$$ g_5(d_a,t_a,d_s,t_s,r_s,t_f,t_w) = E_{used} / E_{total} \le 0.8 $$
	
	
	
	\begin{comment}
		

		1. Introduction
			a. Overview
			b. Goals and objectives
			c. Scope
			d. thesis outline
		2. Literature Review
		3. Methodology
			a. UAV design
				i. UAV geometry 
					a) octocopter, batteries, payload
						i. free body diagram
					b) wingspan calculation
					c) describe drone_geometry module
				ii. mission profile design
				`	a) define mission profile segments 
					b) demonstrate kinematic constraints used
					c) results
				
				iii. battery and propeller selection
				
			c. finite element model
				i. description of the UAV segment model
					a. describe the boundary conditions, materials ect. 
				ii. beam element formulation
					a) timeshenko beams
					b) boundary conditions
					c) mass matrices
				iii. solvers
					a) static 
						i. outputs stress strain
							a) stress and stain plots of th UAV given the payload
					b) dynamic 
						i. inputs 
							a) description of the loads
						ii. outputs amplitude of the response given a force strain, displacement, reactions, ect.
					c) transient
						i. input 
							a) description of the loads
						ii. outputs reactions and displacements over time. Assumes no damping.
					iv. Results 
						
			d. optimization 
				a. objective function: mass calculation
				
				b. design variables
					i. figure showing design variable measurements
					ii. upper and lower bounds 
				
				c. constraints
					i. displacement
						a) describe force
						b) static solver
					ii. stress (treating like black aluminum)
						a) describe force
						b) static solver
					iii. frequency
						a) constraint formula 	
						b) describe forcing function 
						c) dynamic solver
					iv. damage
						a) constraint formula 	
						b) describe forcing function 
						c) transient solver
					v. energy
						a) constraint formula 
						b) show energy consumption calculation
				
				d. results 
	
			a) plots of the design space
			b) show optimized value and compare it to initial
		5. Summary				
		6. Conclusions
		
	
		
	\end{comment}
	
	
		
	
	
	
	
\end{document}
